\documentclass[12pt]{ximera}
\typeout{Start loading xmPreamble.tex}%

% Add here extra macro's that are loaded automatically by all documents of claas 'ximera' or 'xourse' in this repo

%%
%%  Example:
%%
% \newcommand{\R}{\mathbb{R}}
\usepackage{tikz}
\usetikzlibrary{positioning, arrows.meta}
\usepackage{multicol}
\usepackage{enumitem}
\usepackage{caption}
\usepackage{amsmath}
\usepackage{amsthm}

%% NOTE: do NOT add \usepackage to individual activity .tex files.
%% When a xourse inlines an activity it gobbles \documentclass and \input, but a bare
%% \usepackage still executes -- after \begin{document} -- and errors out.
%% All shared packages and macros belong here.
%%
%% ximera.cls already loads: amsmath, amssymb, amsthm, cancel, enumitem, graphicx,
%% hyperref, listings, pgfplots, tikz, url, xcolor. No need to re-load those.

\title{Units and Unit Conversion}
\begin{document}
\begin{abstract}
Converting between units by writing equivalences as fractions, why some conversions
are nonsense, the difference between length, area, and volume units, the metric
system, and converting between measurement systems and currencies.
\end{abstract}
\maketitle

\section*{Units}

Today we will be discussing units and converting between units. In particular, we
will see how they can be useful in doing calculations, be helpful in determining
\emph{what} calculation to do, and how we can use them to show and explain our work.

You will probably want to have your Unit Conversion Sheet available as you work
through today's material.

\subsection*{Converting by writing units as fractions}

Let's begin by revisiting a problem from your pre-class work. Hopefully you were able
to determine the correct answer, but now we will complete it by writing out the units
as fractions in order to convert from one unit to the other.

\begin{example}
How many tablespoons are in a gallon?

We know the following equivalences:
\[
\begin{array}{l}
4 \text{ quarts} = 1 \text{ gallon} \\
4 \text{ cups} = 1 \text{ quart} \\
8 \text{ fl oz} = 1 \text{ cup} \\
2 \text{ Tbs} = 1 \text{ fl oz}
\end{array}
\]

Write each equivalence as a fraction, arranged so the unit you are leaving cancels:
\[
\begin{aligned}
1 \cancel{\text{gal}} \cdot \frac{4 \cancel{\text{quarts}}}{1 \cancel{\text{gal}}}
\cdot \frac{4 \cancel{\text{cups}}}{1 \cancel{\text{quart}}}
\cdot \frac{8 \cancel{\text{fl oz}}}{1 \cancel{\text{cup}}}
\cdot \frac{2 \text{ Tbs}}{1 \cancel{\text{fl oz}}}
&= 4 \cdot 4 \cdot 8 \cdot 2 \text{ Tbs} \\
&= 256 \text{ Tbs}.
\end{aligned}
\]
\end{example}

Now let's try a few other problems.

\begin{problem}
A typical bottle of wine holds about $24.5$ fl oz. How much is this in gallons?

\begin{hint}
Chain the equivalences in the other direction: fl oz $\to$ cups $\to$ quarts $\to$ gallons.
\[
24.5 \cancel{\text{fl oz}} \cdot \frac{1 \cancel{\text{cup}}}{8 \cancel{\text{fl oz}}}
\cdot \frac{1 \cancel{\text{quart}}}{4 \cancel{\text{cups}}}
\cdot \frac{1 \text{ gal}}{4 \cancel{\text{quarts}}}
\]
\end{hint}

A bottle of wine holds about $\answer[tolerance=0.005]{0.19}$ gallons.

\begin{solution}
\[
24.5 \cdot \frac{1}{8} \cdot \frac{1}{4} \cdot \frac{1}{4} \text{ gal}
\approx 0.19 \text{ gal}.
\]
There is about $0.19$ gal in a typical bottle of wine.
\end{solution}
\end{problem}

\begin{problem}
A baby is born and weighs $8$ lbs $4$ oz. Write this as decimal pounds --- that is,
as $8.\underline{\hphantom{00}}$ lbs rather than ``mixing'' pounds and ounces.

The baby weighs $\answer{8.25}$ lbs.

\begin{hint}
There are $16$ oz in $1$ lb, so $4$ oz $= \frac{4}{16}$ lb.
\end{hint}
\end{problem}

\begin{problem}
``How many ounces are in a yard?'' Explain why this question does not make sense.

\begin{freeResponse}
\end{freeResponse}

\begin{solution}
A yard is a unit of \emph{length} and an ounce is a unit of \emph{weight}. They do
not measure the same property, so there is no equivalence between them to convert with.
\end{solution}
\end{problem}

\begin{problem}
``How many feet are in an acre?'' Explain why this question does not make sense.

\begin{freeResponse}
\end{freeResponse}

\begin{solution}
A foot is a unit of \emph{length} (one dimension) and an acre is a unit of
\emph{area} (two dimensions).
\end{solution}
\end{problem}

\subsection*{Length, area, and volume are different properties}

Let's talk a little more about the idea in the question above. It is important to
note that feet, square feet, and cubic feet are all different units --- they are not
measuring the same thing! Feet measure length (one dimension), square feet measure
area (two dimensions), and cubic feet measure volume (three dimensions).

We can convert between feet and inches, or feet and miles. Can we convert between
square feet and square inches, or between cubic feet and cubic miles? To answer this
we have to ask whether the pair of units measures the same property. If the answer is
yes, then yes --- we can convert between them.

\begin{problem}
Do square feet, square inches, and square miles all measure the same property?
\begin{multipleChoice}
\choice[correct]{Yes --- they all measure area.}
\choice{No --- they measure different properties.}
\end{multipleChoice}

Do cubic feet, cubic inches, and cubic miles all measure the same property?
\begin{multipleChoice}
\choice[correct]{Yes --- they all measure volume.}
\choice{No --- they measure different properties.}
\end{multipleChoice}
\end{problem}

In order to convert between units we need some sort of equivalence. Just like knowing
$4$ quarts $= 1$ gallon or $16$ oz $= 1$ pound, these are equivalences we can use to
convert in either direction.

\begin{example}
One \emph{square inch} is the area of a square that is $1$ inch tall and $1$ inch wide.

\begin{center}
\includegraphics[width=3in]{units/units-fig1.jpg}
\end{center}

So a rectangle that is $3$ inches by $4$ inches has area
\[
3 \text{ in} \cdot 4 \text{ in} = 3 \cdot 4 \text{ square inches} = 12 \text{ square inches}.
\]
\end{example}

Let's determine the equivalence between square feet and square inches. A square that
is $1$ foot on each side is the same square as one that is $12$ inches on each side:

\begin{center}
\includegraphics[width=2.5in]{units/units-fig2.jpg}
\qquad
\includegraphics[width=2.5in]{units/units-fig3.jpg}
\end{center}

\begin{problem}
Since the squares are exactly the same size,
\[
1 \text{ sq ft} = \answer{144} \text{ sq in}.
\]

\begin{hint}
$12 \text{ in} \cdot 12 \text{ in} = 12 \cdot 12$ square inches.
\end{hint}

Use the same reasoning to determine the equivalence between square miles and square
feet. Recall that $1$ mile $= 5280$ feet.
\[
1 \text{ sq mi} = \answer{27878400} \text{ sq ft}.
\]

\begin{solution}
$5280 \cdot 5280 = 27{,}878{,}400$ square feet.
\end{solution}
\end{problem}

The same reasoning extends to three dimensions, where a cube $1$ unit on a side gives
us the cubic unit:

\begin{center}
\includegraphics[width=3in]{units/units-fig4.jpg}
\end{center}

\subsection*{The metric system}

Now let's also visit units in the Metric System. You should have had some practice
with these in your pre-class work; let's do a few more together.

\begin{example}
How many decigrams are in a hectogram?
\[
\begin{aligned}
1 \cancel{\text{hg}} \cdot \frac{100 \cancel{\text{g}}}{1 \cancel{\text{hg}}}
\cdot \frac{10 \text{ dg}}{1 \cancel{\text{g}}}
&= 100 \cdot 10 \text{ dg} \\
&= 1000 \text{ dg}.
\end{aligned}
\]
\end{example}

\begin{problem}
$27.3$ kilograms is how many hectograms?
\[
27.3 \text{ kg} = \answer{273} \text{ hg}.
\]

\begin{solution}
\[
27.3 \cancel{\text{kg}} \cdot \frac{1000 \cancel{\text{g}}}{1 \cancel{\text{kg}}}
\cdot \frac{1 \text{ hg}}{100 \cancel{\text{g}}} = 273 \text{ hg}.
\]
\end{solution}
\end{problem}

\begin{problem}
``How many decimeters are in a dekaliter?'' Explain why this question does not make sense.

\begin{freeResponse}
\end{freeResponse}

\begin{solution}
A decimeter measures length; a dekaliter measures volume. Different properties, so
there is no conversion between them.
\end{solution}
\end{problem}

\begin{problem}
Do you notice any patterns in your conversion values?

\begin{freeResponse}
\end{freeResponse}

\begin{solution}
Every metric conversion factor is a power of $10$, so converting only moves the
decimal point. The prefixes tell you how far: kilo ($1000$), hecto ($100$),
deka ($10$), deci ($1/10$), centi ($1/100$), milli ($1/1000$).
\end{solution}
\end{problem}

\subsection*{Converting between systems}

Notice that asking how many feet are in a meter is \emph{not} a nonsense question,
because both measure length. Since inches, feet, miles, and meters all measure
length, there must be some way of converting between the two systems: $5$ feet must
be some number of meters, and $20$ meters must be some number of feet.

To do such conversions we need conversion factors --- equivalences between the
systems. Here is a list of some:
\[
\begin{array}{l}
1 \text{ inch} = 2.54 \text{ cm} \quad\text{and}\quad 1 \text{ cm} = 0.394 \text{ inches} \\
1 \text{ foot} = 0.3048 \text{ m} \quad\text{and}\quad 1 \text{ m} = 3.281 \text{ feet} \\
1 \text{ mile} = 1.609 \text{ km} \quad\text{and}\quad 1 \text{ km} = 0.621 \text{ miles} \\
1 \text{ fl oz} = 29.57 \text{ mL} \quad\text{and}\quad 1 \text{ mL} = 0.338 \text{ fl oz} \\
1 \text{ gal} = 3.785 \text{ L} \quad\text{and}\quad 1 \text{ L} = 0.264 \text{ gal} \\
1 \text{ lbs} = 0.454 \text{ kg} \quad\text{and}\quad 1 \text{ kg} = 2.205 \text{ lbs}
\end{array}
\]

\begin{problem}
Convert your own height into cm (or into inches, if you know your height in cm).
Show the conversion as a chain of fractions.

\begin{freeResponse}
\end{freeResponse}
\end{problem}

\begin{problem}
You have a $43$ lbs (adorable) dog! How much does it weigh in kg?
\[
43 \text{ lbs} \approx \answer[tolerance=0.05]{19.5} \text{ kg}.
\]

\begin{solution}
\[
43 \cancel{\text{lbs}} \cdot \frac{1 \text{ kg}}{2.205 \cancel{\text{lbs}}}
\approx 19.5 \text{ kg}.
\]
\end{solution}
\end{problem}

\begin{problem}
$5$k races are a popular length. How far is a $5$k in miles?
\[
5 \text{ km} \approx \answer[tolerance=0.01]{3.105} \text{ miles}.
\]

\begin{hint}
Use $1 \text{ km} = 0.621$ miles.
\end{hint}
\end{problem}

\begin{problem}
In the U.S.\ we sell gas by the gallon, but most other countries sell and price it by
the liter.

\begin{prompt}
\textbf{(a)} If you have a $15.4$ gal tank, how many liters is that?
\[
15.4 \text{ gal} \approx \answer[tolerance=0.05]{58.3} \text{ L}.
\]
\end{prompt}

\begin{prompt}
\textbf{(b)} Suppose gas here in Columbus is about $\$2.89$ a gallon. How much would
it cost to fill up the $15.4$ gal tank?
\[
\$\answer[tolerance=0.01]{44.51}.
\]
\end{prompt}

\begin{prompt}
\textbf{(c)} Suppose gas in Hamilton, Canada (near Toronto) costs about $\$0.88$
\emph{US dollars} per liter. How much would it cost to fill up the $15.4$ gal tank?
\[
\$\answer[tolerance=0.05]{51.30}.
\]
\end{prompt}

\begin{solution}
\[
15.4 \cancel{\text{gal}} \cdot \frac{3.785 \cancel{\text{L}}}{1 \cancel{\text{gal}}}
\cdot \frac{\$0.88}{1 \cancel{\text{L}}} \approx \$51.30,
\]
so it should cost about $\$51.30$ to fill up the tank.
\end{solution}

\begin{prompt}
\textbf{(d)} Explain why in part (c) we needed the Canadian price in US dollars in
order to make a fair comparison.
\begin{freeResponse}
\end{freeResponse}
\end{prompt}
\end{problem}

\subsection*{Currency is a unit conversion too}

Converting between currencies (such as US and Canadian dollars) is really the same as
a unit conversion --- hence the shared vocabulary, ``converting currency.'' Once you
have an equivalence between the currencies, you can convert one to the other.

\begin{problem}
If $\$1$ CA $= \$0.76$ US, how much does the liter of gas in Hamilton, Canada cost in
Canadian dollars?
\[
\$\answer[tolerance=0.01]{1.16} \text{ CA per liter}.
\]

\begin{solution}
\[
\frac{\$0.88 \text{ US}}{1 \text{ L}} \cdot
\frac{\$1 \text{ CA}}{\$0.76 \text{ US}} \approx \frac{\$1.16 \text{ CA}}{1 \text{ L}}.
\]
\end{solution}
\end{problem}

Looking at the price of gas in the UK (and other parts of Europe), it may seem like
they have very cheap gas --- after all, \$2.89 is a larger number than \pounds 1.26.

\begin{problem}
Calculate how much \pounds 1.26 per liter would be in US dollars per gallon. You
will need to look up the current conversion rate between US dollars and the British
pound.

\begin{freeResponse}
\end{freeResponse}

\begin{solution}
Convert liters to gallons and pounds to dollars in one chain:
\[
\frac{\text{£1.26}}{1 \cancel{\text{L}}} \cdot
\frac{3.785 \cancel{\text{L}}}{1 \text{ gal}} \cdot
\frac{\$\,?\ \text{US}}{\text{£1}}.
\]
At an exchange rate near \pounds 1 to \$1.27 US this is roughly \$6.06 per gallon
--- more than twice the Columbus price, not less.
\end{solution}
\end{problem}

\end{document}
